\chapter{Sĩ Diện Là Món Hàng Đắt}
\markboth{Sĩ Diện Là Món Hàng Đắt}{Pride Is an Expensive Purchase}

\section{Nỗi Sợ Bị Xem Thường}
\begin{truyen}{Tiêu Tiền Để Che Một Cảm Giác}{Spending to Hide a Feeling}
\chuhoa{C}{ó những món người ta mua không phải vì thật sự cần, mà vì sợ ánh mắt của người khác. Một chiếc điện thoại, một bộ quần áo, một bữa ăn sang hơn mức mình đang chịu được.}
\textit{(There are things people buy not because they truly need them, but because they fear the eyes of others. A phone, a set of clothes, a meal more expensive than what they can truly afford.)}

Người thầy cũng từng có những lúc chạnh lòng như vậy. Khi thấy bạn bè đồng nghiệp dùng đồ mới hơn, ăn nói thoải mái hơn về những khoản mua sắm, anh cũng từng thấy mình nhỏ lại. Nhưng cảm giác nhỏ ấy không được chữa lành bằng cách rút ví. Nó chỉ được che đi tạm thời, rồi quay lại còn nặng hơn.
\textit{(The teacher too had moments like that. Seeing colleagues use newer things and speak casually about purchases, he sometimes felt diminished. But that feeling was never healed by opening his wallet. It was only covered temporarily, then returned even heavier.)}

Tiêu tiền vì sĩ diện là kiểu tiêu đáng sợ nhất, bởi nó mua một thứ không thật: cảm giác mình không thua ai. Mà thứ không thật ấy luôn rất chóng tàn.
\textit{(Spending out of pride is among the most dangerous forms of spending, because it purchases something unreal: the feeling that one is not behind others. And that unreal thing fades very quickly.)}
\end{truyen}

\begin{baihoc}
Sĩ diện khiến người ta trả tiền thật cho một cảm giác giả.
\textit{(Pride makes people pay real money for a false feeling.)}
\end{baihoc}

\begin{ghinhoanh}
Pride is costly.

False status drains real money.
\end{ghinhoanh}

\begin{tuhocnhanh}
1. Em có đang mua thứ gì chỉ để ``đỡ ngại'' trước người khác không? \textit{Are you buying something only to avoid embarrassment before others?}

2. Cảm giác nào đang điều khiển việc chi tiêu của em? \textit{Which feeling is controlling your spending?}
\end{tuhocnhanh}
\ngancach

\section{Quản Ninh Không Bỏ Sách}
\begin{truyen}{Nhìn Xe Sang Mà Không Đứng Dậy}{Seeing Splendor Without Rising}
\chuhoa{Q}{uản Ninh và Hoa Hâm cùng ngồi học. Khi có xe ngựa sang trọng đi qua, Hoa Hâm bỏ sách chạy ra xem, còn Quản Ninh vẫn ngồi yên. Câu chuyện nhỏ mà soi được lòng người rất rõ.}
\textit{(Guan Ning and Hua Xin were studying together. When a splendid carriage passed, Hua Xin put down his book and ran out to look, while Guan Ning remained seated. It is a small story, yet it reveals the human heart very clearly.)}

Người ta thường không bị kéo đi bởi chính món đồ bên ngoài, mà bởi ý nghĩ rằng nếu không ngó theo thì mình thiệt, mình kém, mình bỏ lỡ một điều gì lớn. Nhưng người biết mình đang sống vì điều gì thì ít bị kéo lệch đường hơn.
\textit{(People are often drawn away not by the object itself, but by the thought that if they do not look, they are losing, falling behind, or missing something important. Yet those who know what they are living for are less easily diverted.)}

Trong chuyện tiền bạc, giữ được chỗ ngồi của mình giữa sự lấp lánh của người khác là một bản lĩnh rất khó. Nhưng ai không tập điều đó sớm, đời sau này sẽ rất mệt vì cứ phải chạy theo mãi.
\textit{(In financial life, keeping one’s seat amid the glitter of others is a very difficult strength. But those who do not learn it early may later exhaust themselves by chasing endlessly.)}
\end{truyen}

\begin{baihoc}
Người không rời khỏi đường của mình chỉ vì cảnh ngoài thường giữ được tiền và giữ được lòng tốt hơn.
\textit{(Those who do not leave their path merely because of external glitter usually keep both money and inner steadiness more successfully.)}
\end{baihoc}

\begin{ghinhoanh}
Stay in your seat.

Not every glitter deserves your pursuit.
\end{ghinhoanh}

\begin{tuhocnhanh}
1. Điều gì ngoài xã hội làm em mất bình tĩnh nhất? \textit{What in society unsettles you most?}

2. Em có đang chạy theo thứ không thuộc con đường của mình không? \textit{Are you chasing something that does not belong to your path?}
\end{tuhocnhanh}
\ngancach

\section{Một Căn Nhà Yên Không Cần Chứng Minh}
\begin{truyen}{Khi Về Nhà Và Thấy Điều Thật Sự Quan Trọng}{Returning Home to What Truly Matters}
\chuhoa{T}{ối đó, khi người thầy nhìn con đang học bài dưới ánh đèn và nghe tiếng chén bát lạch cạch trong bếp, anh bỗng thấy rõ điều gia đình cần không phải là một hình ảnh sang hơn.}
\textit{(That evening, when the teacher saw his child studying beneath the lamp and heard dishes in the kitchen, he suddenly saw clearly that what the family needed was not a more polished image.)}

Gia đình cần sự vững vàng, cần người cha bình tĩnh, cần những khoản tiền còn lại khi có bệnh tật hay học phí phát sinh. Những thứ ấy không bao giờ được mua bằng sĩ diện. Ngược lại, chúng thường bị sĩ diện âm thầm lấy mất.
\textit{(The family needed stability, a calm father, and money still available when illness or school fees arose. Such things are never purchased through pride. On the contrary, pride often steals them quietly.)}

Từ hôm đó, anh tự nhắc mình: trước khi mua bất cứ thứ gì vì muốn đẹp mặt với người đời, hãy quay về nhìn căn nhà của mình trước. Nếu căn nhà cần bình yên, thì sĩ diện không có quyền đứng đầu.
\textit{(From that day on, he reminded himself: before buying anything to look better in the eyes of others, return first to your own home. If the house needs peace, pride has no right to come first.)}
\end{truyen}

\begin{baihoc}
Gia đình cần sự ổn định nhiều hơn hình ảnh đẹp trước thiên hạ.
\textit{(A family needs stability far more than an impressive image before the world.)}
\end{baihoc}

\begin{ghinhoanh}
Choose peace over image.

Your home matters more than public opinion.
\end{ghinhoanh}

\begin{tuhocnhanh}
1. Nếu bỏ sĩ diện đi, em sẽ giữ được khoản nào cho gia đình? \textit{If you remove pride, what money could you preserve for your family?}

2. Điều gì mới thật sự làm mái nhà của em vững? \textit{What truly makes your home stable?}
\end{tuhocnhanh}

\begin{turanminh}
1. Tôi không được dùng tiền để che cảm giác mặc cảm của mình trước người khác.

2. Cái tôi bị tổn thương không đáng để cả gia đình phải trả giá bằng sự chông chênh cuối tháng.

3. Mỗi khi muốn mua thứ gì chỉ để đỡ ngại với thiên hạ, tôi phải quay về nhìn vợ con, nhìn căn nhà của mình và chọn bình yên thay vì hình ảnh.
\end{turanminh}

\begin{dayhoc}
1. Học sinh rất dễ tiêu tiền vì so sánh. Thầy có thể hỏi: ``Điều em thật sự cần có trùng với điều em đang muốn khoe không?''

2. Từ chương này có thể răn các em rằng giá trị của con người không nằm ở điện thoại, quần áo hay hình thức bên ngoài.

3. Một kết luận dễ nhớ cho lớp là: ``Sĩ diện bắt người ta trả tiền thật cho một cảm giác giả.''
\end{dayhoc}